\chapter{Finite-Difference Methods}\label{ch:qm:finite-difference-methods}

The morning after a grid refinement, a new pricer's gamma near the strike shows a sawtooth. The prices have converged to the cent; their second derivative oscillates from node to node,
positive, negative, positive, with swings twenty times the true gamma, and the hedger's limits flash red for no market reason. Refining the grid again makes the sawtooth twice as large. Nothing
is wrong with the equation or the code in the ordinary sense: the time-stepping scheme is doing exactly what its amplification factor says it will. This chapter discretises the backward equations
of chapter~4 on a grid, analyses consistency, stability and oscillations, and builds the solver the firm uses for one- and two-factor problems.

\section{Discretising a parabolic equation}

By the Feynman--Kac formula, a European price $u(\tau, x)$ in time to maturity $\tau$ and log-price $x = \ln S$ solves $u_\tau = \mathcal Lu$ with $\mathcal Lu = au_{xx} + bu_x - cu$, $a =
\frac12\sigma^2$, $b = r - \frac12\sigma^2$, $c = r$, from the payoff at $\tau = 0$.

\begin{definition}[Finite-difference method, method of lines, non-uniform grid]\label{def:qm:finite-difference-methods:fd}
A \emph{finite-difference method}\index{finite-difference method} replaces the derivatives of a differential equation by difference quotients of the values on a grid of nodes. The \emph{method of
lines}\index{method of lines} discretises space first, turning the equation into a system of ordinary differential equations $\dot U = L_hU$ in time, then applies a time-stepping scheme. On a
\emph{non-uniform grid}\index{non-uniform grid} the spacings $h_i^- = x_i - x_{i-1}$ and $h_i^+ = x_{i+1} - x_i$ differ; the three-point formulas
\[u_{xx} \approx \frac{2u_{i-1}}{h^-(h^- + h^+)} - \frac{2u_i}{h^-h^+} + \frac{2u_{i+1}}{h^+(h^- + h^+)},\]
\[u_x \approx \frac{-h^+u_{i-1}}{h^-(h^- + h^+)} + \frac{(h^+ - h^-)u_i}{h^-h^+} + \frac{h^-u_{i+1}}{h^+(h^- + h^+)}\]
are exact on quadratics.
\end{definition}

\begin{definition}[Truncation error, consistent and stable schemes]\label{def:qm:finite-difference-methods:consistent}
The \emph{truncation error}\index{truncation error} of a scheme is the residual left when the exact solution is substituted into it. A scheme is a \emph{consistent scheme}\index{consistent scheme} if its
truncation error tends to zero with the steps, and a \emph{stable scheme}\index{stable scheme} if its solutions stay bounded, uniformly in the steps, over a fixed time interval.
\end{definition}

\begin{theorem}[Lax equivalence theorem]\label{thm:qm:finite-difference-methods:lax}
For a well-posed linear initial-value problem, a consistent finite-difference scheme converges if and only if it is stable (Lax and Richtmyer, 1956).
\end{theorem}

The theorem divides the work: consistency is Taylor's theorem, and stability is the subject of the next two sections. The firm's operator (\cref{tut:qm:finite-difference-methods:sawtooth}) builds the
tridiagonal $L_h$ on any grid; each implicit step is one solve by the \omterm{def:qm:floating-point-and-numerical-linear-algebra:factor}{tridiagonal matrix algorithm} of chapter~25, $O(n)$ operations. A grid stretched by $x_i = x^* + \alpha\sinh(\xi_i)$ with
uniform $\xi_i$ puts nodes where the payoff bends: with $\alpha = 0.05$ around the strike, it cuts the call's price error by a factor of about 2.9 at equal node count (0.0079 against 0.0225 with 50
nodes, 0.00012 against 0.00035 with 400).

\section{Explicit, implicit and theta schemes}

\begin{definition}[Explicit, implicit, theta and Crank--Nicolson schemes]\label{def:qm:finite-difference-methods:theta}
With time step $\Delta\tau$, the \emph{theta scheme}\index{theta scheme} is $(I - \theta\Delta\tau L_h)U^{n+1} = (I + (1 - \theta)\Delta\tau L_h)U^n$. It is the \emph{explicit scheme}\index{explicit scheme}
(forward Euler) for $\theta = 0$, the \emph{implicit scheme}\index{implicit scheme} (backward Euler) for $\theta = 1$ and the \emph{Crank--Nicolson scheme}\index{Crank--Nicolson scheme} (Crank and Nicolson,
1947) for $\theta = \frac12$; its time error is $O(\Delta\tau)$ except at $\theta = \frac12$, where it is $O(\Delta\tau^2)$.
\end{definition}

\begin{definition}[Von Neumann stability analysis, CFL condition]\label{def:qm:finite-difference-methods:vn}
\emph{Von Neumann stability analysis}\index{von Neumann stability analysis} substitutes a Fourier mode $U_j^n = g^ne^{\mathrm ij\xi}$ into a constant-coefficient scheme and requires $|g(\xi)| \le 1 +
O(\Delta\tau)$ for every frequency. The \emph{CFL condition}\index{CFL condition} (Courant, Friedrichs and Lewy, 1928) is the resulting restriction of the time step by the space step for an \omterm{def:qm:finite-difference-methods:theta}{explicit
scheme}.
\end{definition}

\begin{proposition}[Amplification factors of the theta schemes]\label{prop:qm:finite-difference-methods:amp}
For $u_\tau = au_{xx}$ on a uniform grid, with $\lambda = a\Delta\tau/\Delta x^2$ and $s = \sin^2(\xi/2)$,
\[g(\xi) = \frac{1 - 4(1 - \theta)\lambda s}{1 + 4\theta\lambda s}.\]
The \omterm{def:qm:finite-difference-methods:theta}{explicit scheme} is stable if and only if $\lambda \le \frac12$, that is $\Delta\tau \le \Delta x^2/\sigma^2$; the implicit and \omterm{def:qm:finite-difference-methods:theta}{Crank--Nicolson schemes} are stable for every $\lambda$.
\end{proposition}

\begin{proof}
$U_{j+1} - 2U_j + U_{j-1} = (e^{\mathrm i\xi} - 2 + e^{-\mathrm i\xi})U_j = -4sU_j$, and substituting into the scheme gives $g$. At $\xi = \pi$ the explicit factor is $1 - 4\lambda$, which is at least $-1$ iff
$\lambda \le \frac12$. For $\theta \ge \frac12$ the numerator is at least $-1$ times the denominator for every $\lambda$, so $|g| \le 1$.
\end{proof}

The limit is sharp. With 200 nodes over the grid, the \omterm{def:qm:finite-difference-methods:theta}{explicit scheme} at $\lambda = 0.489$ (409 steps) prices the three-month call at 4.3576 against Black--Scholes 4.3576; at $\lambda = 0.509$ (393
steps) the highest frequency is multiplied by $-1.04$ at each step, and the largest value on the grid is 2\,413.

\section{Stability and oscillations}

Stability in the sense of the Lax theorem is not enough for Greeks. At the highest frequency, the Crank--Nicolson factor is $(1 - 2\lambda)/(1 + 2\lambda)$: bounded by one, but close to $-1$ when
$\lambda$ is large (\cref{fig:qm:finite-difference-methods:amp}). A kinked payoff contains every frequency, and Crank--Nicolson damps its highest ones by almost nothing while flipping their sign at every
step. With 400 nodes and 25 steps for a three-month call ($\lambda = 32$), the price at the strike is within 0.018 of Black--Scholes, and the gamma oscillates between $-0.81$ and $0.55$ around a true
value of 0.040 (\cref{fig:qm:finite-difference-methods:sawtooth}). Refining space and time together doubles $\lambda$ and the sawtooth: the largest gamma error is 0.46, 0.84, 1.69 and 3.38 on 200, 400, 800
and 1\,600 nodes. That is the morning after the refinement.

\begin{omfigure}
\begin{tikzpicture}
\begin{axis}[width=10cm,height=5cm,xlabel={frequency $\xi\Delta x$},ylabel={amplification factor $g$},
  tick label style={font=\small},label style={font=\small},xmin=0,xmax=3.1416,ymin=-1.15,ymax=1.05,grid=major,grid style={gray!25},
  xtick={0,0.7854,1.5708,2.3562,3.1416},xticklabels={$0$,$\pi/4$,$\pi/2$,$3\pi/4$,$\pi$},
  legend columns=2,legend cell align=left,legend style={font=\footnotesize,at={(0.5,-0.3)},anchor=north,draw=none,fill=none,
  /tikz/every even column/.append style={column sep=8pt}},table/col sep=comma]
\addplot[omThm!50,very thick] table[x=xi,y=explicit049]{figdata/methods/27-finite-difference-methods/amplification.csv};
\addplot[omThm,very thick,dashed] table[x=xi,y=explicit051]{figdata/methods/27-finite-difference-methods/amplification.csv};
\addplot[black,very thick] table[x=xi,y=implicit]{figdata/methods/27-finite-difference-methods/amplification.csv};
\addplot[omDef,very thick] table[x=xi,y=cn]{figdata/methods/27-finite-difference-methods/amplification.csv};
\legend{{explicit, $\lambda = 0.49$},{explicit, $\lambda = 0.51$},{implicit, $\lambda = 32$},{Crank--Nicolson, $\lambda = 32$}}
\end{axis}
\end{tikzpicture}

\omcaption{Von Neumann amplification factors of the \omterm{def:qm:finite-difference-methods:theta}{theta schemes} for $u_\tau = au_{xx}$ against the frequency. The \omterm{def:qm:finite-difference-methods:theta}{explicit scheme} leaves $[-1, 1]$ just above $\lambda = \frac12$; at $\lambda = 32$ the \omterm{def:qm:finite-difference-methods:theta}{implicit scheme}
kills the high frequencies and Crank--Nicolson flips their sign with almost no damping. Data: the chapter's tutorial.}\label{fig:qm:finite-difference-methods:amp}
\end{omfigure}

\begin{omfigure}
\begin{tikzpicture}
\begin{axis}[width=11cm,height=5.4cm,xlabel={underlying price $S$},ylabel={gamma},
  tick label style={font=\small},label style={font=\small},xmin=97,xmax=103,ymin=-0.9,ymax=0.65,grid=major,grid style={gray!25},
  legend columns=3,legend cell align=left,legend style={font=\footnotesize,at={(0.5,-0.3)},anchor=north,draw=none,fill=none,
  /tikz/every even column/.append style={column sep=8pt}},table/col sep=comma]
\addplot[omThm!70,thick,mark=*,mark size=1.1pt] table[x=S,y=cn]{figdata/methods/27-finite-difference-methods/sawtooth.csv};
\addplot[black,very thick,dashed] table[x=S,y=exact]{figdata/methods/27-finite-difference-methods/sawtooth.csv};
\addplot[omDef,very thick] table[x=S,y=cn4]{figdata/methods/27-finite-difference-methods/sawtooth.csv};
\legend{Crank--Nicolson,Black--Scholes,{Crank--Nicolson, four implicit half steps}}
\end{axis}
\end{tikzpicture}

\omcaption{Gamma of a three-month at-the-money call from Crank--Nicolson on 400 log-price nodes and 25 time steps ($\lambda = 32$), with and without Rannacher start-up, near the strike (the markers are the grid nodes); the last curve lies on
the Black--Scholes gamma. Data: the chapter's tutorial.}\label{fig:qm:finite-difference-methods:sawtooth}
\end{omfigure}

\begin{definition}[Rannacher time-stepping]\label{def:qm:finite-difference-methods:rannacher}
\emph{Rannacher time-stepping}\index{Rannacher time-stepping} (Rannacher, 1984) replaces the first Crank--Nicolson steps by implicit Euler steps of half the size, which damp the high frequencies of
non-smooth initial data, and continues with Crank--Nicolson.
\end{definition}

Two implicit half steps remove the sawtooth (gamma error $4.7 \times 10^{-4}$ at 400 nodes) but leave a first-order error, halving with each refinement; four restore second order, the error falling
by 4.05 and 4.07 per halving ($1.76 \times 10^{-5}$, $4.35 \times 10^{-6}$, $1.07 \times 10^{-6}$; \cref{fig:qm:finite-difference-methods:conv}). The price tells the same story more quietly: plain
Crank--Nicolson converges only at first order on this problem (errors 0.041, 0.018, 0.0092, 0.0046), with the start-up at second order. Giles and Carter (2006) analyse the mechanism.

\begin{omfigure}
\begin{tikzpicture}
\begin{loglogaxis}[width=10cm,height=5.4cm,xlabel={number of space nodes (time steps $= n/16$)},ylabel={largest gamma error near the strike},
  tick label style={font=\small},label style={font=\small},xmin=150,xmax=2000,ymin=5e-7,ymax=10,grid=major,grid style={gray!25},
  xtick={200,400,800,1600},xticklabels={200,400,800,1\,600},
  legend columns=2,legend cell align=left,legend style={font=\footnotesize,at={(0.5,-0.3)},anchor=north,draw=none,fill=none,
  /tikz/every even column/.append style={column sep=8pt}},table/col sep=comma]
\addplot[omThm,very thick,mark=*,mark size=1.5pt] table[x=n,y=cn]{figdata/methods/27-finite-difference-methods/convergence.csv};
\addplot[black,very thick,mark=square*,mark size=1.5pt] table[x=n,y=implicit]{figdata/methods/27-finite-difference-methods/convergence.csv};
\addplot[omDef!60,very thick,mark=triangle*,mark size=1.8pt] table[x=n,y=cn2]{figdata/methods/27-finite-difference-methods/convergence.csv};
\addplot[omDef,very thick,mark=diamond*,mark size=1.8pt] table[x=n,y=cn4]{figdata/methods/27-finite-difference-methods/convergence.csv};
\legend{Crank--Nicolson,implicit,{two half steps},{four half steps}}
\end{loglogaxis}
\end{tikzpicture}

\omcaption{Largest gamma error of the three-month call for $90 < S < 110$ as space and time are refined together: plain Crank--Nicolson diverges, the \omterm{def:qm:finite-difference-methods:theta}{implicit scheme} and Crank--Nicolson with two implicit
half steps converge at first order, and with four at second order. Data: the chapter's tutorial.}\label{fig:qm:finite-difference-methods:conv}
\end{omfigure}

The payoff's placement on the grid matters as much. A digital paying 1 above the strike, with the strike on a node (where the payoff is set to $\frac12$), converges at first order even with the
start-up (error $9.9 \times 10^{-3}$ on 200 nodes, halving each time) and does not converge at all without it (about 0.1 at every level). With the strike midway between two nodes, which is the same as
averaging the payoff over each node's cell, the error is $6.6 \times 10^{-6}$ on 200 nodes and falls by 4 per halving (\cref{fig:qm:finite-difference-methods:digital}). Pooley, Vetzal and Forsyth (2003)
report the same pattern: averaging the initial data or shifting the grid is not sufficient alone, and restores the expected convergence when combined with a smoothing start.

\begin{omfigure}
\begin{tikzpicture}
\begin{loglogaxis}[width=10cm,height=5cm,xlabel={number of space nodes},ylabel={digital price error},
  tick label style={font=\small},label style={font=\small},xmin=150,xmax=2000,ymin=5e-8,ymax=0.03,grid=major,grid style={gray!25},
  xtick={200,400,800,1600},xticklabels={200,400,800,1\,600},
  legend cell align=left,legend style={font=\footnotesize,at={(0.97,0.5)},anchor=east,fill=white,draw=none},table/col sep=comma]
\addplot[omThm,very thick,mark=*,mark size=1.5pt] table[x=n,y=node]{figdata/methods/27-finite-difference-methods/digital.csv};
\addplot[omDef,very thick,mark=square*,mark size=1.5pt] table[x=n,y=midcell]{figdata/methods/27-finite-difference-methods/digital.csv};
\legend{strike on a node,strike between nodes}
\end{loglogaxis}
\end{tikzpicture}

\omcaption{Price error of a three-month digital at the strike, Crank--Nicolson with four implicit half steps: the strike on a node gives first order, the strike midway between nodes second order.
Data: the chapter's tutorial.}\label{fig:qm:finite-difference-methods:digital}
\end{omfigure}

\section{Boundaries, monotonicity and extrapolation}

\begin{definition}[M-matrix, upwind scheme]\label{def:qm:finite-difference-methods:m}
An \emph{M-matrix}\index{M-matrix} is a real square matrix with nonpositive off-diagonal entries whose inverse exists and is entrywise nonnegative. An \emph{upwind scheme}\index{upwind scheme} discretises
the first derivative by the one-sided difference taken towards the direction the drift comes from, where the central difference would make an off-diagonal coefficient of $L_h$ negative.
\end{definition}

\begin{proposition}[Monotone schemes]\label{prop:qm:finite-difference-methods:m}
If the off-diagonal entries of $L_h$ are nonnegative and $c \ge 0$, the implicit step matrix $I - \Delta\tau L_h$ is an \omterm{def:qm:finite-difference-methods:m}{M-matrix} for every $\Delta\tau$, so the \omterm{def:qm:finite-difference-methods:theta}{implicit scheme} maps nonnegative data to nonnegative
solutions and preserves order: a larger payoff gives a larger price, and no new extremum appears.
\end{proposition}

\begin{proof}
The off-diagonal entries of $I - \Delta\tau L_h$ are nonpositive and each row's diagonal exceeds the sum of the magnitudes of its off-diagonal entries (the rows of $L_h$ sum to $-c \le 0$), so the matrix
is strictly diagonally dominant with a positive diagonal: an \omterm{def:qm:finite-difference-methods:m}{M-matrix}, with a nonnegative inverse.
\end{proof}

Central differences keep the sign pattern while $|b|\Delta x \le 2a$. For a low-volatility, high-rate digital ($\sigma = 2\%$, $r = 5\%$, $\Delta x = 0.02$) the ratio $b\Delta x/(2a)$ is 2.49; the \omterm{def:qm:finite-difference-methods:theta}{implicit
scheme} with central differences then returns a digital worth 1.030 somewhere on the grid, more than the discounted payout of 0.951, and a price that is not monotone in $S$ (total variation 1.23);
upwinding returns 0.951 and a monotone profile, at the price of first-order accuracy in the drift term. At the far ends of the grid the solver imposes either known values (a Dirichlet condition such as
$u = S - Ke^{-r\tau}$ for a deep in-the-money call) or zero gamma, extrapolating linearly; with the grid at five standard deviations the two give prices at the strike that differ by $6 \times 10^{-8}$.

\begin{definition}[Richardson extrapolation]\label{def:qm:finite-difference-methods:rich}
\emph{Richardson extrapolation}\index{Richardson extrapolation} (Richardson, 1911) combines two approximations $u_h$ and $u_{h/2}$ whose error is $Ch^p$ into $u_{h/2} + (u_{h/2} - u_h)/(2^p - 1)$, which
cancels the leading term.
\end{definition}

It works only if the order is what it is assumed to be: with the start-up, the call's price errors on 200 and 400 nodes are $1.8 \times 10^{-3}$ and $4.6 \times 10^{-4}$, and the extrapolated value with $p =
2$ is within $3.2 \times 10^{-7}$. Applied to plain Crank--Nicolson, whose error is first order here, it would remove the wrong term.

\section{Alternating directions}

\begin{definition}[Alternating direction implicit method]\label{def:qm:finite-difference-methods:adi}
An \emph{alternating direction implicit method}\index{alternating direction implicit method} (Peaceman and Rachford, 1955; Douglas and Rachford, 1956) splits a multi-dimensional operator $L = A_0 + A_1 +
A_2$, with $A_1$ and $A_2$ acting along one coordinate each and $A_0$ the mixed-derivative part, and treats one direction implicitly at a time. The Douglas scheme is $Y_0 = U^n + \Delta\tau LU^n$, $Y_j =
Y_{j-1} + \theta\Delta\tau A_j(Y_j - U^n)$ for $j = 1, 2$, $U^{n+1} = Y_2$.
\end{definition}

Each stage solves one tridiagonal system per grid line, so a step costs $O(n^2)$ on an $n \times n$ grid instead of a two-dimensional solve. The mixed derivative is explicit, which is what makes the stability of
the scheme delicate: conditions on $\theta$ for unconditional stability with mixed-derivative terms were analysed by Craig and Sneyd (1988) and by in 't Hout and Welfert (2007), and in 't Hout and
Foulon (2010) compare the Douglas, Craig--Sneyd and Hundsdorfer--Verwer schemes on the Heston equation (stated, not proved here); the firm uses $\theta = \frac12$. On an exchange option $\max(S_1 - S_2, 0)$ with $\sigma_1 = 30\%$, $\sigma_2 = 20\%$ and correlation 0.5, whose price is 10.5243 by Margrabe's formula, the firm's solver
errs by 0.045, 0.026 and 0.014 on grids of $40^2$, $80^2$ and $160^2$ nodes: first order, because the payoff's kink runs along the diagonal between nodes, which the digital above has already shown to be
the usual culprit.

\section{Tutorial: the sawtooth gamma}\label{tut:qm:finite-difference-methods:sawtooth}

\begin{tutorial}
\textbf{Goal.} Price a three-month call on a log-price grid with the three schemes, see the explicit limit, the Crank--Nicolson sawtooth and its removal. \textbf{End state:}
\cref{fig:qm:finite-difference-methods:sawtooth,fig:qm:finite-difference-methods:conv} and the convergence ratios.
\begin{tutsteps}
\item \textbf{The operator.} Three-point coefficients on any grid, with optional upwinding.
\omcode{code/firm/pde/firm_pde.py}{53}{74}{The operator $au_{xx} + bu_x - cu$ on a non-uniform grid (Python).}
\item \textbf{The time step and the start-up}, in the C++20 twin (the Rust twin is line for line the same).
\omcode{code/firm/pde/cpp/firm_pde.hpp}{45}{78}{Theta step with linear boundaries, and Rannacher start-up (C++20).}
\item \textbf{Run} \texttt{explicit\_blowup()}, \texttt{sawtooth()}, \texttt{convergence()}, \texttt{digital\_placement()}, \texttt{stretched()}, \texttt{richardson\_demo()},
\texttt{upwind\_demo()} and \texttt{adi\_exchange()} in \texttt{qm\_fd.py}; then \texttt{fig\_fd.py}.
\end{tutsteps}
\textbf{What to change next.} Price a down-and-out call with the barrier on a node and between nodes; add Rannacher start-up at each discrete monitoring date; try $\theta = 0.51$ instead of the start-up.
\end{tutorial}

\section{Build: the finite-difference solver}\label{bld:qm:finite-difference-methods:pde}

\begin{build}
\bfield{Purpose} One- and two-factor backward equations for the firm's pricers, with Greeks that can be hedged on.
\bfield{Interface} \texttt{uniform\_grid}, \texttt{sinh\_grid}; \texttt{operator(x, a, b, c, upwind)}; \texttt{solve\_1d(x, payoff, a, b, c, tau, n\_steps, theta, rannacher, left, right, upwind)};
\texttt{thomas}; \texttt{richardson}; \texttt{amplification}; \texttt{interp}; \texttt{douglas\_adi\_2d}. C++20 and Rust twins of \texttt{operator\_1d}, \texttt{thomas}, \texttt{theta\_step} and
\texttt{solve\_1d}.
\bfield{Rules} Crank--Nicolson never starts on non-smooth data without implicit half steps; strikes and barriers sit midway between nodes (or the payoff is cell-averaged); the order is measured before
extrapolating; a sign-pattern check flags a lost \omterm{def:qm:finite-difference-methods:m}{M-matrix}.
\bfield{Acceptance tests} \texttt{code/firm/pde/\{tests,cpp,rust\}}: Thomas against a dense solve; the operator exact on quadratics; first and second order measured; the stretched grid beats the
uniform one; Dirichlet boundary; amplification factors against the direct formula; upwinding restores the sign pattern; ADI against Margrabe; C++20 and Rust against the Python values to $10^{-11}$.
\bfield{Stretch} Barriers and American exercise (the penalty method for the \omterm{def:qm:optimal-stopping-and-impulse-control:fbp}{variational inequality} of chapter~10); the Craig--Sneyd and Hundsdorfer--Verwer ADI variants; adaptive grids.
\end{build}

\begin{omsources}
\item J.~Crank and P.~Nicolson, \emph{Proceedings of the Cambridge Philosophical Society} 43, 1947; R.~Courant, K.~Friedrichs and H.~Lewy, \emph{Mathematische Annalen} 100, 1928.
\item P.~D.~Lax and R.~D.~Richtmyer, ``Survey of the stability of linear finite difference equations'', \emph{Communications on Pure and Applied Mathematics} 9, 1956.
\item R.~Rannacher, \emph{Numerische Mathematik} 43, 1984; M.~B.~Giles and R.~Carter, \emph{Journal of Computational Finance} 9(4), 2006; D.~M.~Pooley, K.~R.~Vetzal and P.~A.~Forsyth, ``Convergence remedies
for non-smooth payoffs in option pricing'', \emph{Journal of Computational Finance} 6(4), 2003.
\item D.~W.~Peaceman and H.~H.~Rachford, \emph{Journal of SIAM} 3, 1955; J.~Douglas and H.~H.~Rachford, \emph{Transactions of the AMS} 82, 1956; I.~J.~D.~Craig and A.~D.~Sneyd, \emph{Computers and
Mathematics with Applications} 16, 1988; K.~J.~in 't Hout and B.~D.~Welfert, \emph{Applied Numerical Mathematics} 57, 2007; K.~J.~in 't Hout and S.~Foulon, \emph{International Journal of Numerical Analysis and Modeling} 7, 2010.
\item L.~F.~Richardson, \emph{Philosophical Transactions of the Royal Society A} 210, 1911.
\end{omsources}

\section{Exercises}

\begin{exercise}[$\star$]\label{exo:qm:finite-difference-methods:1}
With $\sigma = 20\%$ and $\Delta x = 0.0025$ in log-price, what is the largest stable explicit time step, and how many steps does a one-year option need?
\end{exercise}

\begin{exercise}[$\star$]\label{exo:qm:finite-difference-methods:2}
Two prices on grids $h$ and $h/2$ are 4.3558 and 4.3572 and the scheme is second order. What is the extrapolated price?
\end{exercise}

\begin{exercise}[$\star$]\label{exo:qm:finite-difference-methods:3}
Show that the three-point second difference on a uniform grid has \omterm{def:qm:finite-difference-methods:consistent}{truncation error} $\frac{h^2}{12}u_{xxxx} + O(h^4)$.
\end{exercise}

\begin{exercise}[$\star\star$]\label{exo:qm:finite-difference-methods:4}
Compute the Crank--Nicolson factor at $\xi = \pi$ for $\lambda = 32$ and the number of steps needed to damp that mode by a factor of 100. Compare with one implicit half step.
\end{exercise}

\begin{exercise}[$\star\star$]\label{exo:qm:finite-difference-methods:5}
For $au_{xx} + bu_x$ with central differences on a uniform grid, find the condition on $\Delta x$ that keeps both off-diagonal coefficients nonnegative.
\end{exercise}

\begin{exercise}[$\star\star$]\label{exo:qm:finite-difference-methods:6}
Why does a Dirichlet condition $u = 0$ at the lower end of a call's grid barely matter, while the same condition at the upper end would ruin the price?
\end{exercise}

\begin{exercise}[$\star\star\star$]\label{exo:qm:finite-difference-methods:7}
\emph{Coding.} Price the three-month digital with the strike on a node and midway between nodes, Crank--Nicolson with four half steps, on 200 to 1\,600 nodes, and give the convergence ratios.
\end{exercise}

\begin{exercise}[$\star\star\star$]\label{exo:qm:finite-difference-methods:8}
\emph{Find the flaw.} ``The price converged to the cent when we halved the grid, so the Greeks we compute by differencing it are right.''
\end{exercise}

\section{Problem: The Sawtooth Gamma}

\begin{problem}[{Weekend problem --- a Greek that gets worse with refinement}]\label{pb:qm:finite-difference-methods:1}
A three-month at-the-money call ($K = 100$, $r = 3\%$, $\sigma = 20\%$) is priced by Crank--Nicolson on a log-price grid of $\pm5$ standard deviations with 400 nodes and 25 time steps.

\medskip\noindent\textbf{Part I --- The schemes.}
\begin{enumerate}
\item What is $\lambda$, and would the \omterm{def:qm:finite-difference-methods:theta}{explicit scheme} be stable?
\item Where exactly does the \omterm{def:qm:finite-difference-methods:theta}{explicit scheme} become unstable, and what happens just beyond?
\item What is the Crank--Nicolson amplification factor of the highest frequency at this $\lambda$?
\item How accurate is the price at the strike, and how accurate is the gamma?
\item What happens to the gamma error when space and time are refined together?
\end{enumerate}

\medskip\noindent\textbf{Part II --- The remedy.}
\begin{enumerate}[resume]
\item What do two implicit half steps give, and at what order?
\item What do four give, and at what order?
\item What order does plain Crank--Nicolson achieve on the price?
\item Why does the digital converge only at first order with the strike on a node?
\item What does moving the strike between nodes give?
\end{enumerate}

\medskip\noindent\textbf{Part III --- Grids and extrapolation.}
\begin{enumerate}[resume]
\item How much does the sinh grid gain at equal node count?
\item What does \omterm{def:qm:finite-difference-methods:rich}{Richardson extrapolation} give with the start-up, and why not without?
\item When do central differences lose the \omterm{def:qm:finite-difference-methods:m}{M-matrix} property, and what does it cost here?
\item How do the two boundary conditions compare?
\item How does Douglas ADI do on the exchange option, and why only first order?
\end{enumerate}

\medskip\noindent\textbf{Part IV --- Judgement.}
\begin{enumerate}[resume]
\item What test would have caught the sawtooth before the hedgers did?
\item Why is a \omterm{def:qm:finite-difference-methods:consistent}{stable scheme} not enough for Greeks?
\item What rules should the firm's pricer enforce by default?
\item State the \emph{named result}: the number of implicit half steps that removes the oscillation and the order it restores, from the error ratio per halving.
\item In one sentence: what does the amplification factor tell a quant that the price does not?
\end{enumerate}
\end{problem}

\section{Interview questions}

\begin{interviewq}[$\star$ \iqroles{researcher, developer}]\label{iq:qm:finite-difference-methods:1}
What is the stability condition of the \omterm{def:qm:finite-difference-methods:theta}{explicit scheme} for the heat equation, and why does it matter for a fine grid?
\end{interviewq}

\begin{interviewq}[$\star\star$ \iqroles{researcher}]\label{iq:qm:finite-difference-methods:2}
Crank--Nicolson is unconditionally stable and second order. Why does it produce oscillating Greeks, and how do you fix it?
\end{interviewq}

\begin{interviewq}[$\star\star$ \iqroles{researcher}]\label{iq:qm:finite-difference-methods:3}
Where do you put the strike and the barrier relative to the grid, and why?
\end{interviewq}

\begin{interviewq}[$\star\star$ \iqroles{developer}]\label{iq:qm:finite-difference-methods:4}
How do you solve an implicit step efficiently in one dimension and in two?
\end{interviewq}

\begin{interviewq}[$\star\star$ \iqroles{researcher, risk}]\label{iq:qm:finite-difference-methods:5}
How would you verify that a finite-difference pricer converges at the order it claims?
\end{interviewq}

\begin{interviewq}[$\star\star\star$ \iqroles{researcher}]\label{iq:qm:finite-difference-methods:6}
Your two-factor solver gives prices that are occasionally negative for a deep out-of-the-money digital. What is happening?
\end{interviewq}
